下面继续在同一冻结前缀上改进 $K$ 的辅助 action 估计。除专门的 logits 结论外，所有向量方差均指协方差 trace，参数固定，动作来自同一 prefix 的独立条件抽样。辅助 action 不等于新增完整 response；对不同 prefix 的 token 套用本节协方差公式会改变目标。

\subsection{一般混合提案的样本上界与输入方差最优性}

记 $\rho=T/q$、$\phi=\rho-1-A$、$c=-\log(1-\eta)/(\kappa\eta)$，这里 $\kappa>0$、$\eta<1$，所以 $c>1$。对任意 $0<\zeta<1$，取 $\mu_\zeta=(1-\zeta)q+\zeta T$、$Z_\zeta=\phi/(1-\zeta+\zeta\rho)$。它对 $K$ 无偏。由 $0\le\phi(x)\le x+c-1$，
\begin{equation}
 0\le Z_\zeta\le C_\zeta,
 \qquad C_\zeta=\max\{(c-1)/(1-\zeta),1/\zeta\}.
\end{equation}
这个上界是逐概率比的精确上确界：$x\to0$ 时 $\phi(x)\to c-1$；$x\to\infty$ 时 $\phi(x)/x\to1$，两个端点分别达到上述两项。第一项随 $\zeta$ 增加、第二项随 $\zeta$ 减少，唯一最小最大值发生在 $\zeta=1/c$，值为 $c$。这解释前节提案的选择，不表示其每个输入的方差最小。

固定输入的方差为
\begin{equation}
 W(\zeta)=\E_q\frac{\phi(\rho)^2}{1+\zeta(\rho-1)}-K^2,
 \qquad W''(\zeta)=2\E_q\frac{\phi(\rho)^2(\rho-1)^2}{[1+\zeta(\rho-1)]^3}\ge0.
\end{equation}
因此可以做一维凸优化选择输入方差最优混合权重；它可以落在端点，也可以不同于 $1/c$。随机 pilot 选权重时，后续样本应与选权重的数据独立，或给出适应性修正。更小的输入方差与更小的统一样本范围是不同目标。
\subsection{把 head--tail 区间变成有界无偏估计器}

令 head 为 $S$、尾部为 $C$，$\delta=\sum_Cq$、$\tau_T=\sum_CT$，假设 $\delta,\tau_T>0$。定义 $r_0=\tau_T/\delta$、$H=\sum_Sq\phi(\rho)$。上一小节的确定性区间可改写为
\begin{align}
 L&=H+\delta\phi(r_0),\\
 U&=H+\tau_T+(c-1)\delta,\\
 w&=U-L=\delta[c+A(r_0)].
 \label{eq:center-tail-interval-lu}
\end{align}
我们可以使随机估计值始终位于这个区间，同时保持无偏。

对尾部概率比定义凸余项
\begin{align}
 g(x)&=\phi(x)-\phi(r_0)-\phi'(r_0)(x-r_0)\\
 &=-A(x)+A(r_0)+A'(r_0)(x-r_0),\\
 u&=A'(r_0)=\frac{r_0^{\kappa-1}}{1-\eta+\eta r_0^\kappa},\qquad
 v=c+A(r_0)-r_0u.
\end{align}
凸性给出 $g\ge0$，且 $A\ge-c$ 给出 $g(x)\le ux+v$。$u>0$；$A$ 凹且 $A(0)=-c$，在 $r_0$ 处的切线于零点的截距 $A(r_0)-r_0A'(r_0)$ 不低于 $-c$，所以 $v\ge0$。尾部中 $\sum_Cq(\rho-r_0)=0$，故
\begin{equation}
 \sum_Cqg(\rho)=K-L,\qquad u\tau_T+v\delta=w.
\end{equation}
当 $w>0$，取仅支持尾部的提案和估计值
\begin{equation}
 \nu(a)=\frac{uT_a+vq_a}{w},\qquad
 Y(a)=\frac{q_ag(\rho_a)}{\nu(a)},\quad a\sim\nu.
\end{equation}
由 $g\le u\rho+v$，有 $0\le Y\le w$。提案可通过条件尾部分布抽样：以概率 $u\tau_T/w$ 从 $T(\cdot\mid C)$ 抽一个 action，否则从 $q(\cdot\mid C)$ 抽一个 action。其响应 rollout 测度仍为 $q$。

\begin{theorem}[区间内的无偏教师中心估计]
\label{thm:center-tail-bounded-estimator}
令 $Y_1,\ldots,Y_M$ 为上述独立尾部样本，$\widehat K=L+M^{-1}\sum_iY_i$。则
\begin{align}
 \E\widehat K&=K,\qquad L\le\widehat K\le U,\\
 \Var\widehat K&\le\frac{(K-L)(U-K)}M\le\frac{w^2}{4M},\\
 \Pr(|\widehat K-K|\ge\epsilon)&\le2\exp(-2M\epsilon^2/w^2).
\end{align}
若 $w=0$ 或尾部为空，直接返回确定的 $K=L=U$。
\end{theorem}

\begin{proof}
尾部余项的恒等式给出 $\E_\nu Y=K-L$，所以均值正确。逐样本范围给出 $Y^2\le wY$，因此 $\Var Y\le w(K-L)-(K-L)^2=(K-L)(U-K)$。独立平均除以 $M$，并用 Hoeffding 得到最后一式。$L\ge0$ 是粗粒化的凸散度，故估计非负。\hfill$\square$
\end{proof}

估计需要能够从条件尾部分布抽样并读取抽中 action 的 $q_a,T_a$。仅有 head 与两种尾部质量时，不能无偏还原 $K$：保持全部已知量不变，让尾部 $T/q$ 恒定得到下界 $L$，将教师尾部质量偏向一个 action 则由严格凸性得到更大的 $K$。仅依赖已知量的随机估计器对两种输入具有相同期望，不能同时无偏。若完整 logits 已经可用，精确求和通常更直接；稀疏缓存模式需把尾部查询成本纳入比较。

无 head 时，$\delta=\tau_T=r_0=1$、$u=1$、$v=c-1$、$L=0$、$w=c$，精确恢复前节全词表包络估计器。扩大嵌套 head 会缩小 $w$，所以最坏方差和置信宽度保证单调改善；不能据此声称实际输入方差也一定单调改善。若允许至少两个未知尾部 action，在固定 head 和尾部质量下可以令尾部 $q$ 质量集中在 $T/q\to0$ 的 action、$T$ 质量集中在极小 $q$ 的另一个 action，使 $K\to U$。任何无偏、非负的残差估计器若有统一样本上界 $b$，都需 $b\ge\sup(K-L)=w$；本构造达到这个信息条件下的统一范围。

还有一个容易实现的备选：直接估计尾部 $\sum_Cq\phi$，使用提案 $(T+(c-1)q)/[\tau_T+(c-1)\delta]$ 在尾部条件化，再加精确 head。固定此包络时，逐步扩大 head 的实际估计方差也不增加：若原提案落入新尾部的概率为 $p_C$，新估计方差为 $p_C^2\Var(Z\mid C)$，原方差由全方差公式至少为 $p_C\Var(Z\mid C)$。但该备选的随机值不保证高于更强的 Jensen 下界 $L$。两种估计分别提供实际方差单调性与更紧的逐样本区间，不混用结论。
